% Laget automatisk av elevark/lag_elevark.py fra sim/videoanalyse.html. Endringer her blir overskrevet.
\documentclass[11pt]{article}
\usepackage{elevark}
\begin{document}
\elevarkhode{Fysikk 1 · Rettlinjet bevegelse}{Videoanalyse uten video}{https://fysikklab.no/sim/videoanalyse.html}{Simuleringen}

Her kan du øve på videoanalyse før du filmer selv. Gå gjennom «videoen» bilde for bilde og klikk på objektet, slik du gjør i Tracker. Send punktene til analysen, og finn fart og akselerasjon.

\section*{Oppgaver}

\begin{oppgaveliste}
\oppgave \fase{test}{Øv}Velg «Ball som slippes». Merk ballen i hvert bilde, og send punktene til analysen. Tilpass en andregradsfunksjon. Hva blir $g$?
\svar
\oppgave \fase{test}{Sammenlign}Last inn eksempeldataene for den samme videoen. Er dine punkter mer eller mindre nøyaktige enn eksempeldataene? Hvordan ser du det på $R^{2}$ og på fartsgrafen?
\svar
\oppgave \fase{test}{Undersøk}Merk bare hvert tredje bilde. Blir resultatet dårligere? Hvorfor kan det likevel være lurt når farten er lav?
\svar
\oppgave \fase{test}{Øv}Velg «Støt mellom to vogner». Merk vogn A og vogn B, og finn fartene før og etter støtet. Er bevegelsesmengden bevart? Bruk tabellen i opplegget Videoanalyse av støt\footnote{\url{fysikklab.no/opplegg/videoanalyse-stot.html}}.
\svar
\oppgave \fase{forklar}{Forklar}Hvorfor blir farten som regnes ut fra nabopunktene, mye mer ujevn enn posisjonen?
\svar
\end{oppgaveliste}

\end{document}
